%% =============================================================
%%  chapters/01_introduction.tex
%% =============================================================
\chapter{Introduction}
\label{ch:introduction}
%% 1-2 pages of content, including references to relevant literature, that explain

\section{Motivation and Relevance}
\label{sec:motivation}
Occupant monitoring is becoming an increasingly important component of modern vehicle safety systems.
Traditionally, vehicle restraint systems were designed around relatively constrained assumptions:
occupants were expected to be seated in approximately upright, forward-facing positions, and 
the driver was expected to remain continuously engaged in the driving task
\cite{mcmurry2018crash}.
These assumptions are becoming less representative as vehicles incorporate higher levels of driving automation, 
more configurable interiors, and a wider range of in-cabin activities. 
In highly automated vehicles, occupants may increasingly recline, rotate their seats, interact with other passengers,
use mobile devices, work, or perform other non-driving-related activities
\cite{mcmurry2018crash}. Consequently,
the vehicle must account not only for whether a seat is occupied, 
but also \emph{who or what occupies the seat and how the occupant is positioned}.




%% why it is important to fuse camera and pressure data, and how it can improve occupant monitoring performance.
%% what are the limitations of single-modality systems.
%% TODO: Expand with concrete regulatory drivers (e.g. Euro NCAP roadmap,
%% GTR requirements) and cite recent OEM occupant monitoring deployments.


\section{Problem Statement}
\label{sec:problem-statement}
Existing occupant classification systems typically rely on a single
sensing modality, which limits their robustness across the wide range of
seating configurations, occupant sizes, and environmental conditions
encountered in real vehicles. Furthermore, collecting sufficiently large
and diverse real-world pressure sensor datasets is costly and raises data
protection concerns, motivating the use of synthetic data generation.
This thesis addresses the problem of designing, training, and evaluating
a multi-task fusion model that jointly estimates occupancy, passenger
class, weight, and pose from synchronised camera and pressure heatmap
inputs.
\cite{sviro} \cite{ticam} \cite{driveandact} \cite{dmd}


%% TODO: Showing the limitations of single-modality systems and the proposed multi-modal fusion approach aim to solve. 
%% 1.Modality imbalance in the literature: published in-cabin datasets and methods are overwhelmingly vision-based, 
%% need for backup due to its limitations (occlusion, lighting, privacy).

%% By iterating EURO NCAP requirements and how the proposed system can help to meet them, 
%% the problem statement can be strengthened.


\section{Research Questions}
\label{sec:research-questions}
The following research questions guide the design, implementation, and
evaluation carried out in this thesis:
" evaluvating suitability of heterogeneous fusion framework for the proposed modalities,comapring its peformance
 to single-modality baselines".

\begin{enumerate}
    \item Fusion Benefits: To what extent does fusing camera and pressure heatmap data
    improve occupancy, passenger classification, weight estimation, and
    pose estimation accuracy compared to single-modality baselines?
    \item Suitability of fusion architecture for these modalities? 
    \item How should synthetic pressure heatmaps be generated and
    validated to ensure they are statistically representative of real
    seat pressure sensor recordings (Sim2Real gap)? \cite{clever2020bodiesrest3dhuman}
    \item Which feature-level fusion strategy provides the best trade-off
    between accuracy and computational cost for a real-time cabin
    monitoring application?
    \item degradation and abstention : how well can it peform under occlusion, lighting changes, faulty pressure sensor data and other adverse conditions, and
     can it abstain from making predictions when one of the modalities is uncertain/low quality input?
\end{enumerate}

\cite{mobihoc} this is a novel deep fusion architecture that can be used to answer the above research questions.


%% TODO: Confirm the final wording of the research questions with the, remove above references.
%% supervisor and ensure they are consistently revisited in Chapter 9.

\section{Objectives and Scope}
\label{sec:objectives-scope}
The primary objective of this thesis is to design and evaluate a
multi-task deep learning architecture that fuses camera images and
pressure heatmaps for cabin occupant state estimation. The scope covers
system requirements definition, data collection , synthetic  pressure heatmaps generation 
generation via human body simulation, a preprocessing and synchronisation
pipeline, the fusion architecture itself, and a quantitative evaluation
including ablation studies.

%% TODO: Add a bullet-point delimitation list of in-scope vs. out-of-scope
%% items once the experimental setup is finalised.

\section{Thesis Structure}
\label{sec:thesis-structure}
Chapter~\ref{ch:background} introduces the background and fundamentals of
cabin monitoring systems, pressure sensing technology, camera systems, and
sensor fusion theory required to follow the remainder of the thesis.
Chapter~\ref{ch:related-work} reviews related work on pressure-based and
camera-based occupant monitoring, multi-modal fusion approaches, and
synthetic data generation, concluding with a gap analysis that positions
this work. Chapter~\ref{ch:system-requirements} defines the use cases,
functional and non-functional requirements, and overall system
architecture. Chapter~\ref{ch:data-acquisition} describes the real data
collection protocol, the synthetic pressure heatmap generation pipeline,
and the Sim2Real validation procedure. Chapter~\ref{ch:preprocessing-sync}
presents the preprocessing and temporal synchronisation pipeline that
prepares camera and pressure data for model training. Chapter
\ref{ch:fusion-architecture} details the multi-task fusion architecture,
loss formulation, and training configuration. Chapter~\ref{ch:results}
reports quantitative results, ablation studies, and Sim2Real transfer
evaluation. Finally, Chapter~\ref{ch:conclusion} summarises the
contributions, answers the research questions, discusses limitations, and
outlines future work.

%% TODO: Update chapter references if chapters are renumbered or merged.
